\documentclass{article}
\usepackage{amsmath,amssymb,amsthm,algorithm,algorithmic}
\usepackage{bm}
\usepackage{hyperref}
\usepackage{enumitem}
\newtheorem{theorem}{Theorem}
\newtheorem{lemma}{Lemma}
\newtheorem{assumption}{Assumption}
\newcommand{\prox}{\operatorname{prox}}
\begin{document}

\section*{Proximal Gradient Method}

\section*{Mathematical Formulation}
We consider the composite convex optimisation problem  

\[
\min_{x\in\mathbb{R}^{d}} \; F(x)\;\stackrel{\text{def}}{=}\; f(x)+g(x),
\]

where  

\begin{itemize}
\item $f:\mathbb{R}^{d}\to\mathbb{R}$ is convex and differentiable with $L$‑Lipschitz continuous gradient,
      \[
      \|\nabla f(x)-\nabla f(y)\|\le L\|x-y\|,\qquad\forall x,y\in\mathbb{R}^{d}.
      \]
\item $g:\mathbb{R}^{d}\to\mathbb{R}\cup\{+\infty\}$ is proper, closed and convex (possibly non‑smooth).
\end{itemize}

For a stepsize $\eta>0$, the \emph{proximal operator} of $g$ is  

\[
\prox_{\eta g}(v)\;\stackrel{\text{def}}{=}\;\arg\min_{u\in\mathbb{R}^{d}}
\Big\{ g(u)+\frac{1}{2\eta}\|u-v\|^{2}\Big\}.
\]

The proximal gradient iteration is  

\[
\boxed{\;x_{k+1}= \prox_{\eta g}\!\big(x_{k}-\eta \nabla f(x_{k})\big)\;},
\qquad k=0,1,2,\dots
\]

\section*{Pseudocode}
\begin{algorithm}[H]
\caption{Proximal Gradient Method}
\begin{algorithmic}[1]
\REQUIRE Initial point $x_{0}\in\mathbb{R}^{d}$, stepsize $\eta\in(0,2/L)$
\FOR{$k=0,1,2,\dots$}
    \STATE Compute gradient $g_{k}\gets\nabla f(x_{k})$
    \STATE Gradient step $y_{k}\gets x_{k}-\eta g_{k}$
    \STATE Proximal step $x_{k+1}\gets\prox_{\eta g}(y_{k})$
    \IF{termination criterion satisfied}
        \STATE \textbf{break}
    \ENDIF
\ENDFOR
\RETURN $x_{k+1}$
\end{algorithmic}
\end{algorithm}

\section*{Dry Run Example}
Minimise $F(w)=f(w)+g(w)$ with  

\[
f(w)=(w-3)^{2},\qquad g\equiv0,
\]
so $\prox_{\eta g}= \operatorname{id}$.  
Take $w_{0}=0$, stepsize $\eta=0.1$.

\begin{tabular}{c|c|c|c}
Iteration $k$ & $w_{k}$ & $\nabla f(w_{k})$ & $w_{k+1}=w_{k}-\eta\nabla f(w_{k})$\\ \hline
0 & $0$ & $2(w_{0}-3)=-6$ & $0-0.1(-6)=0.6$\\
1 & $0.6$ & $2(0.6-3)=-4.8$ & $0.6-0.1(-4.8)=1.08$\\
2 & $1.08$ & $2(1.08-3)=-3.84$ & $1.08-0.1(-3.84)=1.464$\\
3 & $1.464$ & $2(1.464-3)=-3.072$ & $1.464-0.1(-3.072)=1.7712$
\end{tabular}

Objective values $F(w_{k})=(w_{k}-3)^{2}$ decrease from $9$ to $\approx2.60$ after three iterations.

\section*{Assumptions}
\begin{assumption}[Smoothness of $f$]\label{ass:smooth}
$f$ is convex and $L$‑smooth:
\[
f(y)\le f(x)+\langle\nabla f(x),y-x\rangle+\frac{L}{2}\|y-x\|^{2},
\qquad\forall x,y\in\mathbb{R}^{d}.
\]
\end{assumption}
\begin{assumption}[Convexity of $g$]\label{ass:convex}
$g$ is proper, closed and convex.
\end{assumption}
\begin{assumption}[Stepsize]\label{ass:stepsize}
The constant stepsize satisfies $0<\eta<2/L$.
\end{assumption}

\section*{Main Convergence Theorem}
\begin{theorem}[Sub‑linear convergence]\label{thm:main}
Let $\{x_{k}\}$ be generated by the proximal gradient method with a stepsize $\eta\in(0,2/L)$.  
For any optimal solution $x^{\star}\in\arg\min F$, the iterates satisfy  

\[
F(x_{k})-F(x^{\star})\;\le\;
\frac{\|x_{0}-x^{\star}\|^{2}}{2\eta\,k},
\qquad\forall k\ge1.
\]

Consequently $F(x_{k})\to F(x^{\star})$ and the rate is $O(1/k)$.
\end{theorem}

\section*{Theoretical Properties}
\begin{itemize}
\item \textbf{Stability.} The descent inequality (proved below) guarantees that $F(x_{k})$ is non‑increasing for any admissible $\eta$.
\item \textbf{Stepsize sensitivity.} The bound $0<\eta<2/L$ is tight: $\eta=2/L$ yields a non‑expansive mapping but no strict descent; $\eta>2/L$ can break monotonicity.
\item \textbf{Comparison to baselines.}
    \begin{enumerate}[label=\alph*)]
        \item For $g\equiv0$, the method reduces to gradient descent with the classical $O(1/k)$ rate.
        \item Compared with sub‑gradient methods (which achieve $O(1/\sqrt{k})$), exploiting the proximal step yields a faster $O(1/k)$ rate under the same smoothness assumption on $f$.
    \end{enumerate}
\end{itemize}

\section*{Discussion}
The proof hinges on two elementary facts: (i) the optimality condition of the proximal mapping, and (ii) the \emph{non‑expansiveness} of $\prox_{\eta g}$ (Lemma~\ref{lem:nonexp}). By combining these with the $L$‑smoothness of $f$, we obtain a one‑step descent inequality that telescopes across iterations, delivering the $O(1/k)$ bound without any stochasticity or variance terms.

\section*{Preliminaries (Definitions and Lemmas)}
\begin{lemma}[Firm non‑expansiveness of the proximal operator]\label{lem:nonexp}
For any $u,v\in\mathbb{R}^{d}$ and any $\eta>0$,
\[
\big\|\prox_{\eta g}(u)-\prox_{\eta g}(v)\big\|^{2}
\le\langle u-v,\;\prox_{\eta g}(u)-\prox_{\eta g}(v)\rangle .
\]
In particular,
\[
\big\|\prox_{\eta g}(u)-\prox_{\eta g}(v)\big\|
\le\|u-v\|.
\]
\end{lemma}
\begin{proof}
The proximal operator is the resolvent of the maximal monotone operator $\partial g$. Standard results (e.g., \cite{rockafellar1970}) give the firm non‑expansiveness property. ∎
\end{proof}

\section*{Main Proof (Step‑by‑Step Derivation)}
We prove Theorem~\ref{thm:main}.  
Let $x_{k+1}= \prox_{\eta g}(y_{k})$ with $y_{k}=x_{k}-\eta\nabla f(x_{k})$.

\noindent\textbf{[Step 1 (Optimality condition of the prox)}}
By definition of the proximal operator,
\[
\frac{1}{\eta}(y_{k}-x_{k+1})\;\in\;\partial g(x_{k+1}). \tag{1}
\]

\noindent\textbf{[Step 2 (Convexity of $g$)}}
For any $x\in\mathbb{R}^{d}$ and any $s\in\partial g(x_{k+1})$,
\[
g(x)\ge g(x_{k+1})+\langle s,\,x-x_{k+1}\rangle .
\]
Insert $s$ from (1):
\[
g(x)\ge g(x_{k+1})+\frac{1}{\eta}\langle y_{k}-x_{k+1},\,x-x_{k+1}\rangle .
\tag{2}
\]

\noindent\textbf{[Step 3 (Smoothness of $f$)}}
Using $L$‑smoothness (Assumption~\ref{ass:smooth}) with $x=x_{k+1}$ and $y=x_{k}$,
\[
f(x_{k+1})\le f(x_{k})+\langle\nabla f(x_{k}),\,x_{k+1}-x_{k}\rangle
+\frac{L}{2}\|x_{k+1}-x_{k}\|^{2}. \tag{3}
\]

\noindent\textbf{[Step 4 (Combine (2) and (3))}]
Add (2) and (3).  Since $y_{k}=x_{k}-\eta\nabla f(x_{k})$, we have
\[
\langle y_{k}-x_{k+1},\,x-x_{k+1}\rangle
= \langle x_{k}-\eta\nabla f(x_{k})-x_{k+1},\,x-x_{k+1}\rangle .
\]
Thus,
\begin{align*}
F(x) &= f(x)+g(x)\\
&\ge f(x_{k+1})+g(x_{k+1})
   +\frac{1}{\eta}\langle x_{k}-\eta\nabla f(x_{k})-x_{k+1},\,x-x_{k+1}\rangle   \\
&\qquad -\langle\nabla f(x_{k}),\,x_{k+1}-x_{k}\rangle
   -\frac{L}{2}\|x_{k+1}-x_{k}\|^{2}. \tag{4}
\end{align*}

\noindent\textbf{[Step 5 (Rearrange inner products)]}
Observe that
\[
\langle x_{k}-\eta\nabla f(x_{k})-x_{k+1},\,x-x_{k+1}\rangle
-\eta\langle\nabla f(x_{k}),\,x_{k+1}-x_{k}\rangle
= \langle x_{k}-x_{k+1},\,x-x_{k+1}\rangle .
\]
Hence (4) simplifies to
\[
F(x)\ge F(x_{k+1})+\frac{1}{\eta}\langle x_{k}-x_{k+1},\,x-x_{k+1}\rangle
-\frac{L}{2}\|x_{k+1}-x_{k}\|^{2}. \tag{5}
\]

\noindent\textbf{[Step 6 (Choose $x=x^{\star}$)]}
Let $x^{\star}$ be any minimiser of $F$. Substituting $x=x^{\star}$ in (5) gives
\[
F(x_{k+1})-F(x^{\star})
\le \frac{1}{\eta}\langle x_{k}-x_{k+1},\,x_{k+1}-x^{\star}\rangle
+\frac{L}{2}\|x_{k+1}-x_{k}\|^{2}. \tag{6}
\]

\noindent\textbf{[Step 7 (Algebraic identity)}}
For any vectors $a,b$,
\[
2\langle a,b\rangle = \|a\|^{2}+\|b\|^{2}-\|a-b\|^{2}.
\]
Apply with $a=x_{k}-x_{k+1}$ and $b=x_{k+1}-x^{\star}$:
\[
2\langle x_{k}-x_{k+1},\,x_{k+1}-x^{\star}\rangle
= \|x_{k}-x^{\star}\|^{2}-\|x_{k+1}-x^{\star}\|^{2}
-\|x_{k+1}-x_{k}\|^{2}. \tag{7}
\]

\noindent\textbf{[Step 8 (Insert (7) into (6))]}
Dividing (7) by $2\eta$ and using (6):
\[
F(x_{k+1})-F(x^{\star})
\le \frac{1}{2\eta}\Bigl(\|x_{k}-x^{\star}\|^{2}
-\|x_{k+1}-x^{\star}\|^{2}
-\|x_{k+1}-x_{k}\|^{2}\Bigr)
+\frac{L}{2}\|x_{k+1}-x_{k}\|^{2}. \tag{8}
\]

\noindent\textbf{[Step 9 (Collect the $\|x_{k+1}-x_{k}\|^{2}$ terms)]}
Because $0<\eta<2/L$, we have $L-\frac{1}{\eta}>0$ and
\[
\frac{L}{2}-\frac{1}{2\eta}= \frac{L\eta-1}{2\eta}\le 0.
\]
Thus the coefficient of $\|x_{k+1}-x_{k}\|^{2}$ in (8) is non‑positive, and we can drop it to obtain a clean descent inequality:
\[
F(x_{k+1})-F(x^{\star})
\le \frac{1}{2\eta}\Bigl(\|x_{k}-x^{\star}\|^{2}
-\|x_{k+1}-x^{\star}\|^{2}\Bigr). \tag{9}
\]

\noindent\textbf{[Step 10 (Telescoping sum)]}
Sum (9) for $k=0,\dots, N-1$:
\[
\sum_{k=0}^{N-1}\bigl(F(x_{k+1})-F(x^{\star})\bigr)
\le \frac{1}{2\eta}\Bigl(\|x_{0}-x^{\star}\|^{2}
-\|x_{N}-x^{\star}\|^{2}\Bigr)
\le \frac{\|x_{0}-x^{\star}\|^{2}}{2\eta}. \tag{10}
\]

\noindent\textbf{[Step 11 (Bounding the best iterate)]}
Since each term on the left is non‑negative,
\[
N\bigl(F(x_{N})-F(x^{\star})\bigr)
\le \sum_{k=0}^{N-1}\bigl(F(x_{k+1})-F(x^{\star})\bigr)
\stackrel{(10)}{\le}
\frac{\|x_{0}-x^{\star}\|^{2}}{2\eta},
\]
which yields
\[
F(x_{N})-F(x^{\star})\le \frac{\|x_{0}-x^{\star}\|^{2}}{2\eta\,N},
\qquad N\ge1.
\]
This is exactly the claim of Theorem~\ref{thm:main}. ∎

\section*{Rate Derivation (With Constants)}
The constant in the $O(1/k)$ bound is explicitly $\displaystyle \frac{\|x_{0}-x^{\star}\|^{2}}{2\eta}$.  
Choosing the largest admissible stepsize $\eta=1/L$ minimises this constant, giving
\[
F(x_{k})-F(x^{\star})\le \frac{L\|x_{0}-x^{\star}\|^{2}}{2k}.
\]

\section*{Special Cases}
\begin{enumerate}[label=\alph*)]
\item \textbf{Pure gradient descent ($g\equiv0$).} The proximal step becomes the identity, and the proof reduces to the classic gradient‑descent analysis with the same $O(1/k)$ rate.
\item \textbf{Indicator regulariser.} If $g=\iota_{\mathcal{C}}$ is the indicator of a closed convex set $\mathcal{C}$, then $\prox_{\eta g}$ is the Euclidean projection onto $\mathcal{C}$. The algorithm becomes projected gradient descent, and Theorem~\ref{thm:main} recovers the standard convergence guarantee for constrained smooth convex optimisation.
\item \textbf{Strongly convex $F$.} If $F$ is $\mu$‑strongly convex, a refined argument (adding the strong convexity inequality) yields a linear rate
\[
F(x_{k})-F(x^{\star})\le \bigl(1-\eta\mu\bigr)^{k}\bigl(F(x_{0})-F(x^{\star})\bigr).
\]
\end{enumerate}

\bibliographystyle{plain}
\begin{thebibliography}{9}
\bibitem{rockafellar1970}
R.~T. Rockafellar.
\newblock \emph{Convex Analysis}.
\newblock Princeton University Press, 1970.
\end{thebibliography}

\end{document}